\documentclass[12pt, a4paper]{article}

\setlength{\parindent}{0pt}
\setlength{\parskip}{1em}

% Input encoding (for Portuguese accents like ç, ã, é)
\usepackage[utf8]{inputenc}

% Font encoding (standard for modern LaTeX)
\usepackage[T1]{fontenc}

% Language settings for Brazil (e.g., changes "Contents" to "Sumário")
\usepackage[brazil]{babel}

% Page geometry (standard ABNT-like margins)
\usepackage[left=3cm, right=2cm, top=3cm, bottom=2cm]{geometry}

% For clickable links in the PDF (e.g., in references)
\usepackage{hyperref}
\hypersetup{
    colorlinks=true,
    linkcolor=blue,        % Internal links (sections, figures, etc.)
    citecolor=blue,        % Citation links
    urlcolor=blue,         % URL links
    filecolor=blue,        % File links
    menucolor=blue,        % Menu links
    pdftitle={Plano de Trabalho},
    pdfpagemode=FullScreen,
}

% Bibliography support
\usepackage{natbib}

\begin{document}

\title{Proposta de Trabalho - Mestrado}
\author{Cristofer Antoni Souza Costa}
\date{\today}

\maketitle

\textbf{Programa de Pós-Graduação:} Engenharia Mecânica (POSMEC).

\textbf{Linha de Pesquisa:} Mecânica dos Sólidos e Vibrações / Dinâmica de Sistemas Mecânicos.

\textbf{Orientador:} Prof. Dr. Aldemir Aparecido Cavallini Junior, D.Sc.

\section*{Título (provisório)}

\textbf{Diagnóstico de Falhas em Máquinas Rotativas por Meio de Simulações Numéricas e Análise de Espectros de Ordem Superior.}

\section*{Proposta de Plano de Trabalho}
As máquinas rotativas desempenham um papel fundamental em diversos setores industriais, sendo sua operação confiável essencial para a eficiência e segurança dos processos. A detecção precoce de falhas, como trincas ou desalinhamentos, é crucial para evitar paradas não planejadas e reduzir custos de manutenção. Nesse cenário, o monitoramento e a análise de dados de vibração se configuram como as principais ferramentas para este fim \citep{muszynska_rotordynamics_2005}.

Tipicamente, a detecção consiste em identificar a falha por meio de alguma resposta característica do comportamento do eixo em operação permanente ou transiente (aceleração, desaceleração). Nesse sentido, a análise de espectros de ordem superior (Higher-Order Spectra - HOS), incluindo bi-espectro e tri-espectro, oferece uma abordagem promissora para identificar assinaturas não lineares associadas a diferentes tipos de falhas. No entanto, a aplicação prática dessas técnicas muitas vezes enfrenta desafios devido à complexidade dos fenômenos envolvidos e à necessidade de dados experimentais extensivos \citep{sinha_higher_2007}.

Este projeto de mestrado propõe o desenvolvimento de uma metodologia baseada em simulações numéricas para o diagnóstico de falhas em máquinas rotativas, utilizando a análise de HOS. O estudo será conduzido conforme as seguintes etapas:

\begin{enumerate}
\item \textbf{Modelagem Numérica:} Desenvolvimento de modelos de elementos finitos utilizando a ferramenta ROSS \citep{timbo_ross_2020} para simular o comportamento dinâmico de eixos rotativos sob condições normais e com falhas específicas, como trincas e desalinhamentos.

\item \textbf{Análise de Espectros de Ordem Superior:} Aplicação de técnicas de HOS aos dados simulados para identificar padrões característicos que diferenciem as condições normais das falhas modeladas.

\item \textbf{Validação e Avaliação:} Comparação dos resultados obtidos com dados disponíveis na literatura para avaliar a eficácia da metodologia proposta.
\end{enumerate}

Espera-se que este estudo contribua para o aprimoramento das técnicas de diagnóstico de falhas em máquinas rotativas, oferecendo uma alternativa baseada em simulação que possa complementar ou reduzir a dependência de experimentos físicos, alinhando-se às melhores práticas e referências atuais na área.

\bibliographystyle{plain}
\bibliography{references}

\end{document}
